\begin{lemma}[The common-time directions of the barycentric test function]
Let $F$ be a finite-dimensional real inner-product space and let $n\in\N$. Let
$\theta:\{1,\dots,n\}\to\R$ satisfy $\theta_i>0$ for every $i$ and $\sum_i\theta_i=1$,
let $y^0=(y^0_i)_i$ be points of $F$ with $\sum_i\theta_iy^0_i=x_0$, let $t_0>0$ and
$c\in\R$. Let $\varphi$ be a real function on the space-time $\ell^2(F\times\R)$ which is
$C^2$ at $(x_0,t_0)$. On the product space-time $\ell^2\bigl(\ell^2(F^n)\times\R\bigr)$,
write a point as $W=(Y,t)$ with $Y=(y_i)_i$, and let $\Phi$ be a real function there
satisfying, for every such $W$,
\[
  \Phi(Y,t)=\varphi\Bigl(\sum_j\theta_jy_j,\ t\Bigr)+ct
    +\frac{\sum_j\theta_j\bigl|y_j-\sum_k\theta_ky_k\bigr|^2}{2t}.
\]
Then, writing $Y^0=(y^0_i)_i$ and $\nabla_{y_i}\Phi(Y^0,t_0)$ for the $i$-th block of the
spatial part of $\nabla\Phi(Y^0,t_0)$, we have for every $i$
\[
  \theta_i^{-1}\,\nabla_{y_i}\Phi\bigl(Y^0,t_0\bigr)
    =\nabla_x\varphi(x_0,t_0)+\frac{y^0_i-x_0}{t_0},
\]
where $\nabla_x\varphi(x_0,t_0)$ denotes the spatial part of $\nabla\varphi(x_0,t_0)$.
\end{lemma}
\begin{proof}
Write $P=\ell^2(F^n)$ for the spatial factor of the product space-time, and introduce the
continuous linear maps
\[
  \pi_1(W)=Y,\qquad \pi_2(W)=t,\qquad
  X(Y)=\sum_j\theta_jy_j,\qquad
  L(W)=\bigl(X(Y),\,t\bigr),
\]
so that $L$ maps the product space-time to $\ell^2(F\times\R)$, together with, for each
index $j$, the continuous linear map $M_j(W)=y_j-X(Y)$ from the product space-time to
$F$. Each of these is a composition, difference or finite linear combination of
coordinate projections, hence continuous and linear. Set
\[
  N(W)=\sum_j\theta_j\,\bigl|M_j(W)\bigr|^2 ,
\]
so that the hypothesis on $\Phi$ reads
\begin{equation}\label{eq:btd-shape}
  \Phi(W)=\Bigl(\varphi\bigl(L(W)\bigr)+c\,\pi_2(W)\Bigr)+N(W)\cdot\bigl(2\pi_2(W)\bigr)^{-1}.
\end{equation}
Write $W_0=(Y^0,t_0)$. Since $\sum_j\theta_jy^0_j=x_0$ we have
\begin{equation}\label{eq:btd-base}
  L(W_0)=(x_0,t_0),\qquad \pi_2(W_0)=t_0,\qquad M_j(W_0)=y^0_j-x_0 .
\end{equation}
Also, using $\sum_j\theta_j=1$ and $\sum_j\theta_jy^0_j=x_0$,
\begin{equation}\label{eq:btd-off}
  \sum_j\theta_j\bigl(y^0_j-x_0\bigr)=x_0-x_0=0 .
\end{equation}

\emph{Step 1: the derivative of $\Phi$ at $W_0$.}
Because $\varphi$ is $C^2$ at $(x_0,t_0)=L(W_0)$ it is differentiable there, so by the
chain rule for a composition with a continuous linear map, $W\mapsto\varphi(L(W))$ is
differentiable at $W_0$ with derivative $D\varphi(L(W_0))\circ L$. The map
$W\mapsto c\,\pi_2(W)$ is continuous linear, hence differentiable with derivative
$c\,\pi_2$. For each $j$ the map $W\mapsto|M_j(W)|^2$ is differentiable with derivative
$\xi\mapsto2\langle M_j(W_0),M_j\xi\rangle$, since $u\mapsto|u|^2$ has derivative
$v\mapsto2\langle u,v\rangle$ and $M_j$ is continuous linear; multiplying by the constant
$\theta_j$ and summing over the finitely many $j$,
\begin{equation}\label{eq:btd-DN}
  DN(W_0)[\xi]=\sum_j\theta_j\cdot2\bigl\langle M_j(W_0),\,M_j\xi\bigr\rangle .
\end{equation}
Finally $t\mapsto(2t)^{-1}$ is differentiable at $t_0$, because $2t_0\neq0$; composing
with the continuous linear $\pi_2$, the map $W\mapsto(2\pi_2(W))^{-1}$ is differentiable
at $W_0$ with derivative of the form $\mu\,\pi_2$ for some real $\mu$. The product rule
applied to \eqref{eq:btd-shape} therefore gives that $\Phi$ is differentiable at $W_0$
with
\begin{equation}\label{eq:btd-DPhi}
  D\Phi(W_0)[\xi]
   =D\varphi\bigl(L(W_0)\bigr)\bigl[L\xi\bigr]+c\,\pi_2\xi
     +N(W_0)\,\mu\,\pi_2\xi+\frac{DN(W_0)[\xi]}{2t_0}.
\end{equation}

\emph{Step 2: the test direction.}
Fix an index $i$ and a vector $a\in F$, let $e\in F^n$ be the tuple with $e_i=a$ and
$e_j=0$ for $j\neq i$, and put $\xi=(e,0)$ in the product space-time. Then
$\pi_2\xi=0$, and since $\sum_j\theta_je_j=\theta_ia$,
\[
  L\xi=\bigl(\theta_ia,\,0\bigr),\qquad
  M_j\xi=e_j-\theta_ia .
\]

\emph{Step 3: evaluating \eqref{eq:btd-DPhi} on $\xi$.}
The two terms carrying the factor $\pi_2\xi$ vanish. For the first term, by definition of
the gradient and \eqref{eq:btd-base},
\[
  D\varphi\bigl(L(W_0)\bigr)\bigl[L\xi\bigr]
   =\bigl\langle\nabla\varphi(x_0,t_0),\ (\theta_ia,0)\bigr\rangle
   =\bigl\langle\nabla_x\varphi(x_0,t_0),\ \theta_ia\bigr\rangle
   =\theta_i\bigl\langle\nabla_x\varphi(x_0,t_0),\ a\bigr\rangle,
\]
the middle equality because the inner product of the $\ell^2$ product splits as the sum
of the inner products of the two components and the second component of $\xi$ is $0$.
For the last term, by \eqref{eq:btd-DN}, \eqref{eq:btd-base} and Step 2,
\[
  DN(W_0)[\xi]
   =2\sum_j\theta_j\bigl\langle y^0_j-x_0,\ e_j-\theta_ia\bigr\rangle
   =2\sum_j\theta_j\bigl\langle y^0_j-x_0,\ e_j\bigr\rangle
     -2\theta_i\sum_j\theta_j\bigl\langle y^0_j-x_0,\ a\bigr\rangle .
\]
In the first sum only the index $j=i$ contributes, because $e_j=0$ otherwise, and it
contributes $\theta_i\langle y^0_i-x_0,a\rangle$. The second sum equals
$\bigl\langle\sum_j\theta_j(y^0_j-x_0),\,a\bigr\rangle$, which vanishes by
\eqref{eq:btd-off}. Hence $DN(W_0)[\xi]=2\theta_i\langle y^0_i-x_0,a\rangle$ and
\[
  D\Phi(W_0)[\xi]
   =\theta_i\bigl\langle\nabla_x\varphi(x_0,t_0),a\bigr\rangle
     +\frac{\theta_i\bigl\langle y^0_i-x_0,a\bigr\rangle}{t_0}
   =\Bigl\langle\theta_i\Bigl(\nabla_x\varphi(x_0,t_0)+\frac{y^0_i-x_0}{t_0}\Bigr),\ a
     \Bigr\rangle .
\]

\emph{Step 4: reading off the block.}
By definition of the gradient, $D\Phi(W_0)[\xi]=\langle\nabla\Phi(W_0),\xi\rangle$. The
inner product of the product space-time splits as the sum over the two components, and
the time component of $\xi$ is $0$; the inner product of $P=\ell^2(F^n)$ is the sum over
blocks, and only the $i$-th block of $e$ is nonzero. Hence
\[
  \bigl\langle\nabla\Phi(W_0),\xi\bigr\rangle
   =\bigl\langle\nabla_{y_i}\Phi(W_0),\ a\bigr\rangle .
\]
Comparing with Step 3,
\[
  \bigl\langle\nabla_{y_i}\Phi(W_0),\ a\bigr\rangle
   =\Bigl\langle\theta_i\Bigl(\nabla_x\varphi(x_0,t_0)+\frac{y^0_i-x_0}{t_0}\Bigr),\ a
     \Bigr\rangle
\]
for every $a\in F$. As the inner product of $F$ separates points,
\[
  \nabla_{y_i}\Phi(W_0)
    =\theta_i\Bigl(\nabla_x\varphi(x_0,t_0)+\frac{y^0_i-x_0}{t_0}\Bigr),
\]
and dividing by $\theta_i>0$ gives the assertion.
\end{proof}
